\documentclass[UTF8,a4paper]{ctexart}
\usepackage[
  top=2.54cm,
  bottom=2.54cm,
  left=3.17cm,
  right=3.17cm,
  footskip=1.15cm
]{geometry}
\usepackage{amsmath,amssymb}
\usepackage{array}
\usepackage{enumitem}
\usepackage{fancyhdr}
\usepackage{lastpage}

\setmainfont{Times New Roman}
\setCJKmainfont{SimSun}
\setCJKfamilyfont{fs}{FangSong}
\newcommand{\fs}{\CJKfamily{fs}}
\newcommand{\paperfont}{\fs\fontsize{14pt}{23pt}\selectfont}
\newcommand{\checkedbox}{\ensuremath{\rlap{\raisebox{0.1ex}{\checkmark}}\square}}
\newcommand{\emptybox}{\ensuremath{\square}}

\pagestyle{fancy}
\fancyhf{}
\renewcommand{\headrulewidth}{0pt}
\renewcommand{\footrulewidth}{0pt}
\fancyfoot[C]{\fs\fontsize{10.5pt}{12pt}\selectfont 第 \thepage 页 共 \pageref{LastPage} 页}

\setlength{\parindent}{0pt}
\setlength{\parskip}{0pt}
\setlength{\abovedisplayskip}{5pt}
\setlength{\belowdisplayskip}{5pt}
\setlength{\abovedisplayshortskip}{3pt}
\setlength{\belowdisplayshortskip}{3pt}
\setlist[enumerate,1]{
  label=\arabic*.,
  leftmargin=1.55em,
  labelsep=0.35em,
  itemsep=0pt,
  topsep=0pt,
  parsep=0pt,
  partopsep=0pt
}
\setlist[enumerate,2]{
  label=(\arabic*),
  leftmargin=2.5em,
  labelsep=0.25em,
  itemsep=0pt,
  topsep=0pt,
  parsep=0pt,
  partopsep=0pt
}
\setlist[itemize]{
  label=\textbullet,
  leftmargin=2.4em,
  itemsep=0pt,
  topsep=0pt,
  parsep=0pt,
  partopsep=0pt
}

\newcommand{\examrowheight}{\rule{0pt}{27pt}}
\newcolumntype{L}[1]{>{\raggedright\arraybackslash}p{#1}}

\begin{document}
\paperfont

\begin{center}
{\songti\bfseries\fontsize{16pt}{20pt}\selectfont 复旦大学研究生课程考试试卷}\\[13pt]
{\fontsize{14pt}{18pt}\selectfont （2025  - 2026 学年第 2 学期）}
\end{center}

\vspace{8pt}

{\fontsize{14pt}{17pt}\selectfont
\setlength{\tabcolsep}{0pt}
\renewcommand{\arraystretch}{1.02}
\begin{center}
\begin{tabular}{@{}L{2.45cm}L{5.55cm}L{2.45cm}L{4.15cm}@{}}
\examrowheight 授课学期： & 春季 & 考试日期： & 2026-7-1 \\
\examrowheight 课程名称： & 神经网络的模型与应用 & 课程代码： & MATH40054 \\
\examrowheight 开课院系： & 数学科学学院 & 试卷类型： & \checkedbox A 卷\quad \emptybox B 卷 \\
\examrowheight 考试形式： & \multicolumn{3}{L{12.15cm}}{\emptybox 开卷\quad \emptybox 闭卷\quad \emptybox 课程论文\quad \checkedbox 其他} \\
\multicolumn{4}{@{}L{14.6cm}@{}}{\examrowheight 姓名：\quad 盖括\hspace{1.65cm} 学号：\quad 25110180015\hspace{1.30cm} 成绩：} \\
\end{tabular}
\end{center}
}

\vspace{24pt}

\begin{center}
（共7题，每题20分，选5道题完成即可）
\end{center}

\vspace{28pt}

\begin{enumerate}
\item 对于如下过程描述的LIF模型
\[
\frac{dV}{dt}=-g_LV+g_{\mathrm{syn}}I_{\mathrm{syn}}
\]
\[
\tau\frac{dI_{\mathrm{syn}}}{dt}=-I_{\mathrm{syn}}+\mu+\sigma\frac{dW}{dt}
\]
这里W是一个标准的Wiener过程，膜电位达到阈值$V_{th}$产生动过电位放电，存在复位周期$\tau_{ref}$。$I_{syn}$描述（一个）突触电流，$\tau$是其时间常数。$\mu$、$\sigma$、$g_{syn}$、$\tau$、$g_L$皆为常数。对此动作电位发放点稳态过程的Inter-Spike-Interval的随机变量记为$T$：
\par (1) 给出$T$的期望和方差的表达式，推导该神经元动作电位发放的更新函数(renewal function）的期望和方差表达式；
\par (2) 给定任意的期望$m$和方差$\sigma^2$，如何构造参数$(g_{syn},\mu,\sigma)$（给定
\newpage
其他常数）,使得$T$的期望和方差分别等于或者近似$m$和$\sigma^2$，并通过数值模拟来验证你的结果。

\item 对于一个ordinary renewal process（$X_1,X_2,\cdots,X_n,\cdots$），假设已知其更新间隔$X_n$各阶统计量，请推导该神经元动作电位发放的更新函数(renewal function）的n阶统计量的表达式。

\item 考虑如下兴奋-抑制神经网络模型
\[
\tau_E\frac{dv_E}{dt}=-v_E+\left[M_{EE}v_E-M_{EI}v_I+h_E\right]^+
\]
\[
\tau_I\frac{dv_I}{dt}=-v_I+\left[M_{IE}v_E-M_{II}v_I+h_I\right]^+
\]
$[z]^+=\max\{z,0\}$, $v_{E,I}$分别表示兴奋性/抑制性神经元群的发放率，$M_{ab}$分别是对应链接权重，a,b=E,I, $h_{E,I}$分别是兴奋性/抑制性神经元群接收到的外部刺激。分别给出该系统对应参数$\tau_I$和$h_E$的Hopf分叉分析，并给以实际数值模拟进行验证。

\item 多神经元对刺激方向的响应模型如下。假设每个神经元的条件发放率表示如下
\[
f_a(s)=\exp\left(-\frac{(s-s_a)^2}{2\sigma_a^2}\right)
\]
其中a=1，…,N是神经元编号，s是受到外部刺激的方向角度（$[0,180^\circ]$），$s_a$是神经元a的偏好方向，$\sigma_a>0$是神经元调谐函数的参数。假设
\newpage
\par {\Large$\bullet$}\quad 神经元动作电位发放之间是统计独立的；
\par {\Large$\bullet$}\quad 神经元动作电位发放是一个泊松过程。
\par 假设贝叶斯损失为平方差：
\[
loss=(s-\hat{s})^2
\]
$\hat{s}$是估计方向角度。

通过这N个神经元的动作电位发放时间点，建立对外部刺激方向s的最小方差的无偏(MVUB)估计，并计算其平方误差。可给出解析解，或者近似解，并进行数值模拟验证。

\item 分离附件中音频数据中的源信号并评估质量，并且比较至少两种算法。

\item 构建神经网络对MNIST数据集进行0-9手写数字的分类，并分别利用随机梯度下降（SGD）和自然梯度（NG）进行训练，并进行对比。

\item 利用Reinforcement Learning算法求解迷宫最短路径搜索问题。\par 迷宫图请见附件(maze.jpg；黄色是出发位置)。
\end{enumerate}

要求：任意给出中点位置，规划出最短路径或者告知不能到达。最好给出交互式的界面。

\end{document}
